\chapter{Dodecaphase incidence and the Witt--Leech--Monster corridor}
\label{chap:bandera-excepcional-acumulativa}
\index{double projection!common flag}
\index{Witt!design}
\index{Monster}

The \cref{chap:paley-codigo-witt} has explicitly constructed the matrix
\(A_W\), the code \(C_W\) and the system \(W_{12}\). On those coordinates,
the double projection is formalized by two evaluation maps which,
in the declared gauge, produce the same polarized Witt flag through distinct
operations. This calibrated flag also determines a marked origin. The
construction reaches the design \(W_{12}\), its automorphism group
\(M_{12}\) and, through a second branch that does not pass through \(W_{24}\), an
even unimodular rootless 3-neighbor. Classical classification identifies the
latter with the Leech lattice; the extension to the moonshine vertex
operator algebra and the Monster group is obtained through the corresponding classical
theorems.

\section{Dodecaphase state and incidence supports}

The dodecaphase vector is
\[
K=(234,543,140,729,659,824,621,058,914,794,146,601).
\]
The four words of the frame are
\[
w_\pi=010211,
\quad
w_e=201101,
\quad
w_\varphi=121200,
\quad
w_{\rm APP}=022020.
\]
In the symmetric Paley--Witt gauge the encoder used is
\[
c(w)=(w,wA_W)\in\mathbb F_3^{12}.
\]
Their supports are
\[
\begin{aligned}
P_\pi&=\{2,4,5,6,7,8,9,10,11\},\\
H_e&=\{1,3,4,6,9,10\},\\
H_\varphi&=\{1,2,3,4,7,9\},\\
H_{\rm APP}&=\{2,3,5,10,11,12\}.
\end{aligned}
\]
\(P_\pi\) has nine points. The other three supports are hexads. This
distinction must be preserved: the twelve hexads contained in \(P_\pi\) have
the lines of \(AG(2,3)\) as complementary triples.

\section{Coincidence of the threshold supports}

Archimedean evaluation extracts from the vector \(K\) the threshold support
\[
\beta_{\rm arq}(K)=\{i:K_i\ge729\}
=\{4,6,9,10\}.
\]
Exceptional evaluation calculates the intersection
\[
\beta_{\rm exc}(w_e,w_\pi)
=\operatorname{supp}c(w_e)\cap\operatorname{supp}c(w_\pi).
\]
Evaluation produces the first exact equality conditional on these data:
\begin{equation}
\boxed{
\beta_{\rm arq}(K)
=\beta_{\rm exc}(w_e,w_\pi)
=B_K=\{4,6,9,10\}.}
\label{eq:bandera-cara-alta}
\end{equation}
The two maps perform distinct operations on declared
inputs: the first uses thresholds of the vector \(K\), whereas the second
uses intersections of code supports. The threshold \(729\), the words and
the Paley gauge are part of the calibration. In particular, this
equality is not an independent validation of \(K\): all integer thresholds
between \(660\) and \(729\), inclusive, select the same support
\(B_K\).

\section{Coincidence of the negative supports}

Before carrying, the Archimedean cochain is
\[
s^{(0)}=P+E-\Phi-K,
\]
and its twelve components are
\[
(7,297,353,-430,-715,-1199,-3,286,-894,-528,380,663).
\]
Its negative support is
\[
\eta_{\rm arq}(s^{(0)})
=\{4,5,6,7,9,10\}.
\]
By the exceptional route, we fix as calibration selector the hexad \(H\)
satisfying
\[
B_K\subset H\subset P_\pi
\]
and
\[
(|H\cap H_e|,|H\cap H_\varphi|,|H\cap H_{\rm APP}|)=(4,3,2).
\]
Enumeration of the design produces a unique solution for that selector:
\[
\eta_{\rm exc}
=\{4,5,6,7,9,10\}
=H^-_\alpha.
\]
Therefore,
\begin{equation}
\boxed{
\eta_{\rm arq}(s^{(0)})
=\eta_{\rm exc}(B_K,P_\pi,H_e,H_\varphi,H_{\rm APP})
=H^-_\alpha.}
\label{eq:bandera-hexada-negativa}
\end{equation}

Equations~\eqref{eq:bandera-cara-alta} and
\eqref{eq:bandera-hexada-negativa} are condensed into the flag
\[
\mathcal I_*^{\rm core}=(B_K\subset H^-_\alpha).
\]
This is the calibrated compatibility flag of the double projection. The marked frame
\[
\widetilde{\mathcal I}_*
=(\mathcal I_*^{\rm core};P_\pi,H_e,H_\varphi,H_{\rm APP})
\]
preserves the data needed to construct the exceptional branch. Under a
simultaneous permutation \(g\in M_{12}\), the entire construction is transported
equivariantly; the invariant is its isomorphism class, not the isolated numerical
names.

\section{Reconstruction of the base point from the incidence frame}

The origin need not remain an external selector if the complete ordered
frame is preserved. We define
\[
o(\mathcal F)
=(H_e\cap H_\varphi)\setminus(H^-\cup H_{\rm APP}),
\qquad
\mathcal F=(H^-,H_e,H_\varphi,H_{\rm APP}).
\]

\begin{theorem}[Origin determined by incidences]
For the canonical frame,
\[
o(\mathcal F_*)=\{1\}.
\]
Moreover, the rule is equivariant: \(o(g\mathcal F)=g\,o(\mathcal F)\).
\end{theorem}

\begin{proof}
We have
\[
H_e\cap H_\varphi=\{1,3,4,9\},
\]
and
\[
(H^-_\alpha\cup H_{\rm APP})\cap\{1,3,4,9\}
=\{3,4,9\}.
\]
The difference is \(\{1\}\). Equivariance follows because intersection,
union and difference commute with the action of any permutation.
\end{proof}

The passage from this incidence origin to a vector of a lattice
\(A_2^{12}\) requires fixing a metric realization of the twelve coordinates.
The \cref{sec:vecino-leech-ejecutable} fixes that realization, shows
that the origin selects one of the twelve radial minimizers and certifies
the conditions of the 3-neighbor. Incidence determines the origin; the Gram
matrix and the positive chamber then determine the vector.

\section{Incidence of the ternary design}

The Paley chain
\[
C_P\longmapsto A_W\longmapsto C_W\longmapsto W_{12}
\]
is fixed by the matrices and proofs of the
\cref{chap:paley-codigo-witt}. In particular,
\(A_W^2=-I_6\), \(C_W=[12,6,6]_3\), and the 132 projective supports of
weight six form \(W_{12}=S(5,6,12)\). Passing from a word to its support
identifies exactly the pair \(\{c,-c\}\); the incidences used in
what follows therefore belong to the projective design and not to individual signed
words.

\section{Hexads, octads and the cardinality identity 28}

The binary design \(W_{24}=S(5,8,24)\) has octads, each with
\(\binom82=28\) pairs. The
\cref{p12-83-c17-s1:thm:dodecada-w12} constructs
\(W_{12}\) as the residual design of a dodecad \(D\) of the extended binary
Golay code; the \cref{thm:elevacion-unica-hexada} proves that, once
the alignment \(\ell\) is fixed, each hexad has a unique lift
\(O_H\) to an octad. The step \(6\to8\) is therefore an incidence
morphism relative to the binary datum \((D,\ell)\), and not an inference based
on cardinalities.

On the other hand, the microfiber of \(\pi\) contains 1008 routes and a nonad has
36 pairs, so that
\[
1008=36\cdot28.
\]
The identity \(1008=36\cdot28\) is a numerical check of the catalogue. The
incidence morphism is defined independently by means of
\(H\mapsto\iota_D(\ell(H))\); the cardinality identity does not enter into
its definition.

\section{Star involutions and generation of the Mathieu group}

For each face \(B\in\binom{[12]}4\), there are four hexads of \(W_{12}\)
containing \(B\). The four residues \(H\setminus B\) are disjoint pairs
partitioning \([12]\setminus B\). The involution \(\sigma_B\) fixes \(B\)
and exchanges the endpoints of those four pairs. There are
\[
\binom{12}{4}=495
\]
star involutions.

Exhaustive enumeration, for the declared matrix and words, verifies
\[
\sigma_B(W_{12})=W_{12}
\]
for all of them and
\[
\left|\langle\sigma_B:B\in\tbinom{[12]}4\rangle\right|=95040.
\]
Since \(\Aut(W_{12})\cong M_{12}\), the 495 stars generate \(M_{12}\).
The HMT face is
\[
B_K=\{4,6,9,10\},
\qquad
\sigma_{B_K}=(1\ 3)(2\ 11)(5\ 7)(8\ 12).
\]
A single star generates only \(C_2\). The global combinatorial action is
expressed through the representation of the free group of star generators
\[
\rho_\star:F_{495}\longrightarrow\Aut(W_{12}),
\qquad
\rho_\star(\gamma_B)=\sigma_B,
\qquad
\operatorname{im}\rho_\star=M_{12},
\]
where \(F_{495}\) is the free group on the 495 formal stars. The global
group is therefore the image of that family of involutions and not that of a
single star. Identifying these generators with loops in the state space
would additionally require constructing the corresponding transport; that
identification does not enter into the group calculation.

\section{Two exceptional extensions of the ternary code}

Two branches of different types can be continued from the ternary code:
\[
\begin{aligned}
C_W&\longrightarrow W_{12}
\xrightarrow[\ (D,\ell)\ ]{\ \iota_D\ }W_{24},\\
C_W&\xrightarrow[\phi]{\text{explicit gluing}}N(A_2^{12})
\xrightarrow[v_\alpha]{\text{exact 3-neighbor}}\Lambda_{24}
\longrightarrow V^\natural .
\end{aligned}
\]
The first branch continues through the residual design of the dodecad and its
unique lift; the second continues through ternary gluing and an index-three
neighbor. The two constructions have distinct codomains. For the
second, let \(A_2=\mathbb Z\alpha_1\oplus\mathbb Z\alpha_2\) with Gram
matrix
\[
G_{A_2}=\begin{pmatrix}2&-1\\-1&2\end{pmatrix},
\qquad \rho=\alpha_1+\alpha_2,
\]
and take as representatives of \(A_2^*/A_2\cong\mathbb F_3\)
\[
\varpi_0=0,\qquad
\varpi_1=\frac{2\alpha_1+\alpha_2}{3},\qquad
\varpi_2=\frac{\alpha_1+2\alpha_2}{3}.
\]
The gluing realization is the coordinate-by-coordinate map
\[
\phi(c_1,\ldots,c_{12})=(\varpi_{c_1},\ldots,\varpi_{c_{12}}),
\]
which forms
\[
N(C_W)=\bigcup_{c\in C_W}(A_2^{12}+\phi(c)).
\]
The linearity of \(C_W\) makes the union of cosets a lattice; its
self-duality gives the integrality of the gluing, and divisibility by three of
all weights gives its evenness. Moreover, \(C_W\) has dimension six and
\[
\det N(C_W)=\frac{3^{12}}{|C_W|^2}
=\frac{3^{12}}{(3^6)^2}=1.
\]
Minimum distance six implies that a nonzero gluing class has norm
at least \(6(2/3)=4\). Therefore the roots of \(N(C_W)\) are exactly those
of \(A_2^{12}\). This is the explicit construction of the Niemeier of type
\(A_2^{12}\) used below.

\section{Radial minimization and explicit construction of the 3-neighbor}
\label{sec:vecino-leech-ejecutable}

In the positive chamber we consider radial vectors
\[
v(a)=\sum_{i=1}^{12}a_i\rho_i,
\qquad a_i=1+3b_i,\quad b_i\ge0.
\]
The condition \(v(a)^2\equiv0\pmod{18}\) is equivalent to
\(\sum_i b_i\equiv1\pmod3\). Its minimum is
\[
v(a)^2=2\sum_i(1+3b_i)^2=54,
\]
and is attained exactly twelve times: one coefficient is four and the remaining eleven
are one. The recovered origin \(o_*=1\) selects
\[
v_\alpha=4\rho_1+\rho_2+\cdots+\rho_{12}.
\]
Minimality belongs to the declared positive radial cone. In the entire
residual quotient there exists, for example, the representative
\((\alpha_1-2\alpha_2,\rho,\ldots,\rho)\), congruent to
\((\rho,\ldots,\rho)\) modulo three and of norm \(14+11\cdot2=36\). The
direct calculation preserves it as a witness that prevents extending the domain.

We define
\[
N_{\alpha,0}=\{x\in N(C_W):
\langle x,v_\alpha\rangle\equiv0\pmod3\},
\qquad
N_\alpha=N_{\alpha,0}+\mathbb Z\frac{v_\alpha}{3}.
\]

\begin{theorem}[Leech 3-neighbor]
The preceding construction satisfies
\begin{align*}
\mathrm{(I1)}\quad&v_\alpha\in A_2^{12}\subset N(C_W),
&\langle N(C_W),v_\alpha\rangle&\not\subset3\mathbb Z,\\
\mathrm{(I2)}\quad&v_\alpha^2=54\equiv0\pmod{18},\\
\mathrm{(I3)}\quad&v_\alpha/3\in
N_{\alpha,0}^*\setminus N(C_W),
&3(v_\alpha/3)&\in N_{\alpha,0}.
\end{align*}
The lattice \(N_\alpha\) is even, unimodular, of rank twenty-four and
rootless. Consequently,
\[
N_\alpha\cong\Lambda_{24}.
\]
\end{theorem}

\begin{proof}
For any root of a factor \(A_2\), the pairing with
\(a_i\rho_i\) is congruent to \(1\) or \(2\) modulo three. The character of
(I1) is therefore surjective, and none of the old roots belongs to
\(N_{\alpha,0}\). The equality in (I2) is the preceding radial calculation. If
\(x\in N_{\alpha,0}\), then
\(\langle x,v_\alpha/3\rangle\in\mathbb Z\), which proves the first
inclusion in (I3). If \(v_\alpha/3\) belonged to \(N(C_W)\), the
integrality of the Niemeier would force the character of (I1) to be zero. The
class has order three and \(v_\alpha\in N_{\alpha,0}\) because
\(v_\alpha^2\equiv0\pmod3\).

The kernel has index three, hence
\[
\det N_{\alpha,0}=3^2\det N(C_W)=9,
\qquad
\det N_\alpha=9/3^2=1.
\]
For \(x\in N_{\alpha,0}\) and \(m\in\mathbb Z\),
\[
\left(x+m\frac{v_\alpha}{3}\right)^2
=x^2+2m\frac{\langle x,v_\alpha\rangle}{3}+6m^2
\in2\mathbb Z,
\]
so the neighbor is even. It remains to exclude new roots. Each coordinate of
the two nontrivial cosets belongs to one of the six classes
\[
A_2+\varpi_j\pm\rho/3,\qquad j\in\{0,1,2\},
\]
and the minimum of each class is \(2/9\). Indeed, writing a coordinate
as \((u/3,v/3)\), its norm is
\(2(u^2-uv+v^2)/9\); the six residues are nonzero and each admits a
representative with \(u^2-uv+v^2=1\). Every
new norm is at least \(12(2/9)=8/3>2\). The neighbor has no roots. The
classification of even unimodular lattices of rank twenty-four
identifies it with Leech.
\end{proof}

Enumeration of the six local minima and the twelve radial minimizers
reconstructs the \(729\) words of \(C_W\), the Gram matrix of the gluing and
its evenness; exact arithmetic checks I1--I3, the determinants and the bound
\(8/3\).

\section{Corridor from Leech to the Moonshine algebra}
\label{sec:corredor-hmt-moonshine}

Define
\[
 \Lambda_{\rm HMT}:=N_\alpha\cong\Lambda_{24},
 \qquad
 V^\natural_{\rm HMT}
 :=V_{\Lambda_{\rm HMT}}^+
 \oplus V_{\Lambda_{\rm HMT}}^{T,+}.
\]
No second lattice is introduced here: the
Frenkel--Lepowsky--Meurman construction is applied to the neighbor produced in the preceding
subsection.

\begin{theorem}[HMT--Moonshine corridor]
\label{thm:corredor-hmt-moonshine}
With the dodecad, alignment, origin and exceptional chamber of the
preceding construction fixed, we have
\[
 V^\natural_{\rm HMT}\cong V^\natural,
 \qquad
 \Aut(V^\natural_{\rm HMT})\cong\mathbb M,
\]
and
\[
 \operatorname{Tr}_{V^\natural_{\rm HMT}}q^{L_0-1}
 =J_{\rm Mon}(\tau)
 =j(\tau)-744
 =q^{-1}+196884q+21493760q^2+\cdots .
\]
Consequently, the chain
\[
 \mathrm{HMT}\longrightarrow N_\alpha\cong\Lambda_{24}
 \longrightarrow V^\natural_{\rm HMT}
 \longrightarrow\mathbb M
 \longrightarrow\text{\rm Moonshine}
\]
is a result proved by composition.
\end{theorem}

\begin{proof}
The 3-neighbor theorem proves, from the declared data, that
\(N_\alpha\) is even, unimodular, of rank twenty-four and rootless. Lattice
classification identifies it with \(\Lambda_{24}\). Applied to
this output, the Frenkel--Lepowsky--Meurman construction produces the
orbifold
\(V_{\Lambda_{\rm HMT}}^+\oplus V_{\Lambda_{\rm HMT}}^{T,+}\);
the classical theorems identify its automorphism group with
\(\mathbb M\) and its graded character with the normalized Hauptmodul
\(J_{\rm Mon}=j-744\)
\cite{FLM1988,ConwayNorton1979,Borcherds1992}. Each arrow therefore delivers
an object of the type required by the next.\qedhere
\end{proof}

\begin{remark}[The trivial component is not a fit]
The grading satisfies
\[
 V^\natural_{\rm HMT}
 =\bigoplus_{n\geq0}(V^\natural_{\rm HMT})_n,
 \qquad
 (V^\natural_{\rm HMT})_0=\CC\mathbf1,
 \qquad
 (V^\natural_{\rm HMT})_1=0.
\]
The term \(q^{-1}\) thus records the vacuum line. In degree two,
\[
 \dim(V^\natural_{\rm HMT})_2
 =196884
 =1+196883
 =196560+324
 =54(5\cdot729+1).
\]
The equality \(196884=1+196883\) separates the trivial representation and the smallest
nontrivial irreducible representation of the Monster. The other two equalities
are exact integer decompositions; they neither replace this graded structure
nor are used to select the corridor retrospectively.
\end{remark}

The finite stage up to \(N_\alpha\) starts from a declared dodecad, alignment,
origin and chamber. The calculation proves that \(N_\alpha\) is even,
unimodular and rootless. The classical classification and orbifold theorems
are applied to that output: the orbifold is \(V^\natural\), its automorphism
group is \(\mathbb M\) and its graded character is \(J_{\rm Mon}\).
The composition thus preserves the provenance of each arrow without reducing the
corridor to a coincidence of coefficients.

\section{Quotient space of 243 syndromes}

Puncturing one coordinate of \(C_W\) produces
\[
G_{11}=[11,6,5]_3.
\]
The ternary ball of radius two has
\[
1+22+220=243
\]
elements. Since
\[
|G_{11}|\cdot243=3^6\cdot3^5=3^{11},
\]
and the minimum distance is five, these balls are disjoint and cover all of
\(\mathbb F_3^{11}\). Each vector admits a unique decomposition
\[
x=c+e,
\qquad
c\in G_{11},
\quad
\operatorname{wt}(e)\le2.
\]
The linear quotient is
\[
\mathbb F_3^{11}/G_{11}\cong\mathbb F_3^5,
\qquad
|\mathbb F_3^{11}/G_{11}|=243.
\]
Thus, a quotient space of 243 syndromes is obtained in an ambient space of eleven
coordinates and quotient dimension five. The catalogue of 243 words \(w_6\)
and this quotient have the same cardinality, but are objects of different types;
the direct encoder sends code words to the zero syndrome. The proved double projection
uses the common flag, not an unstructured cardinality
bijection.

\section{Finite compatibility of the arithmetic and exceptional evaluations}

\begin{theorem}[Finite compatibility of the two projections]
Consider the vector \(K\) and the signed dodecaphase state constructed in
the \cref{chap:alpha}; the four catalogue words defined in the
\cref{chap:regiones}; and the matrix \(A_W\) constructed in the
\cref{chap:paley-codigo-witt}. The construction fixes a single calibration
\(c_{\rm exc}\in\mathscr C_{\rm exc}\) and transports it with the state; this
calibration types the coordinates of the exceptional codomain and does not act as
a generating input. The following statements hold:
\begin{enumerate}
  \item Archimedean evaluation obtains \(\alpha_{12}\) through dodecaphase
  carrying;
  \item the Archimedean and exceptional support maps produce the same flag
  \(B_K\subset H^-_\alpha\);
  \item the marked frame recovers the origin \(o_*=1\);
  \item the matrix \(A_W\) constructs \(W_{12}\), and the 495 star
  involutions generate a group of order \(95040\), identified with \(M_{12}\);
  \item the explicit discriminant realization constructs \(N(A_2^{12})\),
  and the origin, together with the positive chamber, selects \(v_\alpha\); its
  3-neighbor is even, unimodular and rootless.
\end{enumerate}
The support identities and the origin rule are equivariant under a
simultaneous permutation of the twelve coordinates. The calibration transforms the
incidence data simultaneously and does not alter their isomorphism class.
\end{theorem}

\begin{proof}
The first statement is the dodecaphase reconstruction of the
\cref{chap:alpha}. Equations
\eqref{eq:bandera-cara-alta} and \eqref{eq:bandera-hexada-negativa} prove
the second by direct evaluation of the two maps. The set-theoretic calculation
in the preceding section gives \(o_*=1\). Finally, the enumerator, the 132
weight-six supports, the 495 involutions and the group order, together with the
3-neighbor theorem, prove the last statement.
\end{proof}

\begin{corollary}[Compatibility with the HMT--Moonshine corridor]
The finite data of the exceptional projection construct the same
\(N_\alpha\) that enters into the
\cref{thm:corredor-hmt-moonshine}; therefore, finite evaluation is
compatible with \(V^\natural_{\rm HMT}\), its \(\mathbb M\) action and its
character \(J_{\rm Mon}\).
\end{corollary}


The finite flag is obtained by two distinct operations applied to
coordinates of the same enriched state, and the origin is derived after
fixing that flag. This proves a typed internal identity and preserves the
common provenance without promoting calibration to a causal input. The
lattice construction preserves the type distinction. The composition from
the data of \(\alpha\) to \(\Lambda_{24}\) factors through the code, the
gluing \(A_2^{12}\), the marked origin, the radial chamber and the 3-neighbor. The
branch \(W_{12}\to W_{24}\) continues separately and is determined for every
choice of a dodecad \(D\) and an incidence isomorphism
\(\ell:W_{12}^{\rm HMT}\to\mathcal H(D)\). The proved equivariance means
that those choices coordinate the representation without altering the generated
isomorphism class or the closure of the double projection.
